\section{PolyCone: search for nonlinear inequality certificates}

\subsection{Certificate language and its soundness}

Let $p,g_1,\ldots,g_m,f_1,\ldots,f_r\in\Q[x_1,\ldots,x_n]$.
Assume the current context proves $g_i(\vec x)\geq0$ and
$f_j(\vec x)=0$. A certificate for $p(\vec x)\geq0$ is the exact polynomial
identity
\begin{equation}\label{eq:cone}
 p=\sum_{k=1}^{s}\lambda_k q_k^2\prod_{i=1}^{m}g_i^{e_{ki}}
   +\sum_{j=1}^{r}h_j f_j,
 \qquad \lambda_k\in\Q_{\geq0},\quad e_{ki}\in\N,
\end{equation}
where $q_k,h_j$ are rational polynomials. Repetition of a factor is permitted;
the empty product is 1. The equality multipliers $h_j$ may have either sign.

\begin{proposition}[Cone-certificate soundness]\label{prop:cone}
If identity \eqref{eq:cone} holds and all of its sign and equality hypotheses
hold at a real valuation, then $p$ is nonnegative at that valuation.
\end{proposition}
\begin{proof}
Every square is nonnegative. Products of the nonnegative $g_i$, the square,
and $\lambda_k$ are nonnegative, so their sum is nonnegative. Each $h_jf_j$
vanishes at the valuation. Evaluating the polynomial identity gives the claim.
\end{proof}

This proof is small compared with the search for the identity. The implemented
checker independently reconstructs both sums using rational sparse-polynomial
arithmetic, checks factor indices and coefficient signs, and compares the
result with $p$ exactly. It does not consult the LP solver or a numerical
residual threshold. Its input hypotheses remain hypotheses: a Lean adapter
must supply their actual proofs in the current scope.

\subsection{Finite dictionary construction}

For degree cap $D$, generate a bounded collection of columns
\[
 c_k=q_k^2\prod_i g_i^{e_{ki}},\qquad \deg c_k\leq D.
\]
Begin with monomial $q_k$ of degree at most $D/2$. Add binomial candidates
suggested by the target: unit sums and differences, exact rational square-root
ratios of positive even-degree terms, and completing-square expressions
inferred from cross coefficients. If
\[
 a u^2+cuv+bv^2
\]
appears, with $a>0$, try $q=u+(c/(2a))v$; when $b>0$, also try the opposite
orientation. This is a proposal mechanism, not a rule asserting that the
remaining polynomial is nonnegative.

The completing-square proposal matters in practice. The initial heuristic
recognized $(x-2/3)^2$, but a positive extra constant obscured its coefficient
ratio in $(x-2/3)^2+1/97$. Adding the cross-coefficient construction recovered
the intended square. A regression test now covers this case. The unsuccessful
initial search returned \code{unknown}; it did not accept an approximate proof.

Form products of known nonnegative constraints subject to the same degree cap.
Include columns $m f_j$ for monomials $m$ of degree at most $D-\deg f_j$; their
multipliers are unrestricted rational variables. Deduplicate identical columns
and retain their construction recipes. Avoid expanding the global Cartesian
product when most constraints concern unrelated variables: partition by
variable support and request cross-component products only when the target
connects those components.

The shipped implementation uses a simple bounded enumeration and a limit of
4,000 generators. It is a research implementation, not an optimized sparse
SOS package. A production implementation needs a separate bound on enumeration
work, monomial count, coefficient bit length, and external-solver time; a cap
on successfully admitted columns alone is insufficient as a denial-of-service
or worst-case complexity control.

\subsection{Search becomes linear algebra with an untrusted oracle}

Expand the columns and target in a common monomial basis. The certificate
search is
\begin{equation}\label{eq:lp}
 A\lambda+Bh=b,\qquad \lambda\geq0,
\end{equation}
where the $h$ coordinates are free. The Python prototype calls SciPy's
\code{linprog(method="highs")} to propose a solution \cite{scipy-lp}.
It minimizes the sum of cone coefficients as a simple heuristic. A production
objective should also penalize large polynomial terms, expensive guard proofs,
and predicted replay cost; minimizing coefficient sum is not proof-size
minimization.

Floating-point coefficients are converted to rational candidates with bounded
denominators and checked exactly. If that fails, solve the proposed active
support with exact SymPy linear algebra, set any remaining free parameters to
zero, and check the result again. A negative reconstructed coefficient, an
inconsistent support, or a failed identity causes rejection. No tolerance is
used at the acceptance boundary.

\begin{Code}
polycone(target, inequalities, equalities, D):
  columns := degree-bounded products and proposed squares
  add ideal columns with unrestricted coefficients
  numerical_candidate := solve_linear_program(columns, target)
  if no candidate: return unknown
  rational_candidate := reconstruct(numerical_candidate)
  if exact_check(rational_candidate): return certificate
  support_candidate := exact_solve(numerical_support)
  if exact_check(support_candidate): return certificate
  return unknown
\end{Code}

The choice of zero values for free reconstruction parameters is incomplete:
a different rational point on the same support could have nonnegative
coefficients. Similarly, numerical infeasibility does not establish
nonnegativity failure and may not even establish infeasibility of the exact
finite problem. These are performance limitations only because every
successful result passes an independent exact check.

\subsection{Worked certificates}

Several tested cases explain the range of the certificate language.
For the two-dimensional Cauchy--Schwarz inequality,
\[
 (a^2+b^2)(c^2+d^2)-(ac+bd)^2=(ad-bc)^2.
\]
For variance in three variables,
\[
 x^2+y^2+z^2-xy-yz-zx
 =\tfrac12(x-y)^2+\tfrac12(y-z)^2+\tfrac12(z-x)^2.
\]
Both follow from squares without order assumptions on the variables.
For $x,y\in[0,1]$,
\[
 xy(1-x)(1-y)\geq0
\]
uses four known nonnegative factors. The identity has degree 4 and is rejected
by the prototype's degree-2 search before LP construction.

Equality hypotheses contribute differently. Under $y=x$,
\[
 y^2-2x+1=(x-1)^2+(x+y)(y-x).
\]
The second term need not be nonnegative; it vanishes because of the equality.
Treating all certificate multipliers as nonnegative would unnecessarily lose
this useful algebraic flexibility.

\subsection{Theory cooperation through generated inequalities}

The plugin should return useful consequences before it can prove the final
goal. Suppose $l_x\leq x\leq u_x$, $l_y\leq y\leq u_y$, and a shared atom
$z$ is certified equal to $xy$. The four products of interval distances yield
\begin{align*}
 z&\geq l_x y+l_y x-l_xl_y, &
 z&\geq u_x y+u_y x-u_xu_y,\\
 z&\leq u_x y+l_y x-u_xl_y, &
 z&\leq l_x y+u_y x-l_xu_y.
\end{align*}
For instance, the first follows from $(x-l_x)(y-l_y)\geq0$.
These inequalities are linear in the shared atoms $x,y,z$ and can strengthen
an existing linear arithmetic solver. Their derivations remain polynomial
certificates with explicit bound dependencies. The proposal is to generate
such cuts on demand, not to replace the linear engine.

For natural numbers and integers, use their actual semantics. Natural
subtraction is truncated; machine arithmetic may wrap; division and remainder
have domain-specific rules. Reify into an ordered field only after proving the
necessary cast and operation correspondence. A real-arithmetic relaxation can
supply a sufficient certificate for an integer theorem after a checked
embedding, but it cannot justify an integer witness or an integer-specific cut
without an additional argument.

Strict inequalities require separate bookkeeping. A positive rational margin
in a certificate is enough to prove strict positivity; nonnegativity alone is
not. Rational-expression normalization requires denominator conditions. These
conditions are obligations on the fact graph, not informal side notes in an
external certificate.

\subsection{Limits and a concrete SOS escalation}

The number of monomials of total degree at most $D$ in $n$ variables is
$\binom{n+D}{D}$. Even before proof reconstruction, indiscriminate degree and
variable growth can make the coefficient problem too large. Apply cheap
symbolic factorization and variable separation first, then expand a bounded
basis in increments, recording which expansion helped.

A finite binomial-square dictionary is strictly more restrictive than arbitrary
sum-of-squares search. An optional escalation can introduce a Gram matrix
$Q\succeq0$ with $p=v^{\mathsf T}Qv$ or its constrained analogue, solve for $Q$
using a semidefinite optimizer, and reconstruct an \emph{exact} rational
factorization
\[
 Q=L\,\operatorname{diag}(d_1,\ldots,d_t)L^{\mathsf T},\qquad d_i\geq0.
\]
The resulting weighted squares fit the same certificate checker. A numerical
claim of positive semidefiniteness is not accepted; an exact factorization and
identity are required. Singular matrices and rational reconstruction can make
this escalation fail. No SDP solver is implemented in this package.

The tested Motzkin polynomial
\[
 M(x,y)=x^4y^2+x^2y^4+1-3x^2y^2
\]
is nonnegative by arithmetic--geometric mean applied to the three nonnegative
terms $x^4y^2$, $x^2y^4$, and 1. The finite unconstrained square dictionary does
not certify it. Its status is therefore \code{unknown}, despite truth. No
bounded cone in this report is advertised as a complete real-algebraic decision
procedure. Rational multipliers, richer positivity theorems, or other proof
methods require explicit new certificate forms and soundness lemmas rather
than a more generous floating-point tolerance.

\section{Bernstein certificates on rational boxes}

\subsection{A complementary proof language}

A local box constraint can make a positivity proof easier than a global
square decomposition. Map $x\in[a,b]$, $a<b$, to $t=(x-a)/(b-a)\in[0,1]$.
For degree $d$ write the polynomial exactly as
\[
 p(a+(b-a)t)=\sum_{k=0}^{d} b_k B_{k,d}(t),\qquad
 B_{k,d}(t)=\binom{d}{k}t^k(1-t)^{d-k}.
\]
If every $b_k\geq0$, each summand is nonnegative on the interval. Products of
these bases give a certificate for a multivariate box.

For power coefficients $a_j$ on $[0,1]$, the conversion formula is
\[
 b_k=\sum_{j=0}^{k}a_j\frac{\binom{k}{j}}{\binom{d}{j}}.
\]
To verify it, use
$t^j=\sum_{k=j}^{d}\binom{k}{j}\binom{d}{j}^{-1}B_{k,d}(t)$,
which follows by expanding the binomial theorem for the remaining degree
$d-j$. The multivariate conversion is its tensor product.

The prototype discovery routine computes the coefficients after an affine
change of variables. Its checker takes a deliberately different route:
it reconstructs every basis term directly in the original coordinates,
\[
 \binom{d}{k}\frac{(x-a)^k(b-x)^{d-k}}{(b-a)^d},
\]
and checks the resulting polynomial identity exactly. Both routines share the
sparse rational arithmetic layer, so their implementation independence is not
complete. Differential tests against SymPy address that shared layer but do
not constitute formal verification.

\subsection{Subdivision is a proof tree, not sampling}

If some coefficients are negative, bisect a box coordinate and try both
children. A successful certificate is a complete binary tree of boxes, with
nonnegative-coefficient leaves. The checker verifies the root box, each cut's
location, the exact two child boxes, and every leaf identity. A gap in the
cover invalidates the entire certificate even if every retained leaf looks
correct. Testing finitely many points is never an acceptance rule.

For
\[
 p(x)=x^2-x+\tfrac13\quad\text{on }[0,1],
\]
the degree-2 coefficients are $(1/3,-1/6,1/3)$, so the initial box does not
certify positivity. On $[0,1/2]$ they become $(1/3,1/12,1/12)$, and on
$[1/2,1]$ they become $(1/12,1/12,1/3)$. Two leaves prove the entire interval.
This example is included in the tests and measurements.

\subsection{Why true statements can remain unresolved}

The polynomial $(x-1/3)^2$ is nonnegative, yet midpoint subdivision never makes
$1/3$ an endpoint of an interval with dyadic endpoints. On the leaf containing
that point in its interior, all Bernstein basis functions are strictly
positive there. If every coefficient were nonnegative and their sum were zero
at that interior point, every coefficient would have to vanish. The polynomial
would then vanish identically on the leaf, which it does not. Thus this
particular subdivision strategy cannot finish, not merely because the tested
depth cap is too small.

This is a useful division of labor: the cone routine can recognize the square,
while a naive box routine cannot. Conversely, box constraints can certify
polynomials that a small global dictionary misses. Route a problem according
to observed structure, preserve both alternatives, and report the actual
certificate family that succeeded.

A production splitter can use rational root isolation or symbolic factorization
to expose zero sets before subdivision. That is a concrete improvement, but
it requires certified root information and complete domain decomposition.
Adaptive numerical guesses about root locations remain proposals until those
obligations are discharged.
